\chapter{Deriving the Innovation Gating Parameter}
\label{chap:deriving-alpha}

This chapter provides a self-contained derivation of the alpha parameter 
used in the innovation gating mechanism, as introduced in \cref{eq:alpha_lambda} of \cref{chap:theory-framework}. 
The derivation follows closely the work of \citeonline{Fortaleza:2022}, from where the gating mechanism was inspired.

Let $Y \in \mathbb{R}$ be a random variable representing the output of a given process in time, and $Y_I$ the observation of $Y$.
It is assumed that the observation is affected by a Gaussian noise $N \sim \mathcal{N}(0, \sigma_n^2)$, such that
%
\begin{equation}
    Y_I = Y + N.
\end{equation}

Let $y_I$ be a realization of $Y_I$.
Denoting by $y_M$ the prediction of $y$ from a given model, 
$y_O$ is defined to be the adjusted prediction after the observation $y_I$ is taken into account.
For the sake of simplicity, $y_O$ is given by the following exponential moving average (EMA) relationship:
%
\begin{equation}
    \begin{cases}
        y_O[0] = y_I[0], & k = 0, \\
        y_O[k] = \alpha[k] y_I[k] + (1 - \alpha[k]) y_M[k-1], & k > 0
    \end{cases}
\end{equation}
%
where $\alpha[k] \in [0, 1]$ is the smoothing factor of the EMA at time step index $k$.

Parameter $\alpha[k]$ represents the statistical confidence that 
the filter output reflects the same direction of the process output at sample $k$. 
More precisely, $\alpha[k]$ is given by the following expression:
%
\begin{align}
    \alpha[k] = \ \
    &P(\mathrm{sign}(Y_M[k-1] - Y_I[k]) =    \mathrm{sign}(\mu_{Y_M[k-1]} - \mu_{Y_I[k]})) \nonumber\\ -
    &P(\mathrm{sign}(Y_M[k-1] - Y_I[k]) \neq \mathrm{sign}(\mu_{Y_M[k-1]} - \mu_{Y_I[k]})),
\label{eq:alpha}
\end{align}
%
where $Y_M[k-1]$ is an observation of $Y_M$ with mean $\mu_{Y_M[k-1]}$
and $Y_I$ an observation of $Y_I$ with mean $\mu_{Y_I[k]}$. 
Here $\mathrm{sign}(x)$ represents the sign function, which is defined as:
%
\begin{equation}
    \mathrm{sign}(x) =
    \begin{cases}
        1, & x > 0, \\
        0, & x = 0, \\
        -1, & x < 0.
    \end{cases}
\end{equation}

\Cref{eq:alpha} can be rewritten as follows:
%
\begin{equation}
    \alpha[k] = 
    2 P(\mathrm{sign}(Y_M[k-1] - Y_I[k]) = \mathrm{sign}(\mu_{Y_M[k-1]} - \mu_{Y_I[k]})) 
    - 1.
\label{eq:alpha2}
\end{equation}

Defining a radom variable $Y_{MI}$ whose observation at time $k$ have mean and variance given by
%
\begin{align}
    \mu_{Y_{MI}[k]} &= |\mu_{Y_M[k-1]} - \mu_{Y_I[k]}|, \\
    \sigma^2_{Y_{MI}[k]} &= \sigma^2_{Y_M[k-1]} + \sigma^2_{Y_I[k]} - 2 \mathrm{Cov}(Y_M[k-1], Y_I[k]),
\end{align}
%
where $\mathrm{Cov}(Y_M[k-1], Y_I[k])$ is the covariance between $Y_M[k-1]$ and $Y_I[k]$, 
\cref{eq:alpha2} can be expressed as follows:
%
\begin{equation}
    \alpha[k] = 2 P(Y_{MI}[k] > 0) - 1.
\label{eq:alpha3}
\end{equation}

In order to find a practical implementation of $\alpha[k]$, \cref{eq:alpha3} is analyzed under the following
technical assumptions.

\begin{assumption}
    The random variable $Y_{MI}[k]$ can be approximated by a normal distribution 
    with mean $\mu_{Y_{MI}[k]}$ and variance $\sigma^2_{Y_{MI}[k]}$.
\label{assumption:normal}
\end{assumption}

\begin{assumption}
    Both high frequency and random behavior of the measurement noise implies that 
    the observations $Y_M[k - 1]$ and $Y_I[k]$ can be considered approximately statistically independent of each other. 
    Mathematically, $\mathrm{Cov} (Y_M[k-1], Y_I[k]) \approx 0$.
\label{assumption:independence}
\end{assumption}

\begin{assumption}
    The variance of the model's prediction is much smaller than the variance of the measured observation.
    Mathematically, $\sigma^2_{Y_M[k-1]} \ll \sigma^2_{Y_I[k]}$.
\label{assumption:variances}
\end{assumption}

\begin{assumption}
    By considering an additive zero-mean measurement noise, the samples $y_M[k-1]$ and $y_I[k]$ 
    can be used to estimate $\mu_{Y_M[k-1]}$ and $\mu_{Y_I[k]}$, respectively.
\label{assumption:means}
\end{assumption}

From \cref{assumption:normal}, \cref{eq:alpha3} can be rewritten as:
%
\begin{align}
    \alpha[k] 
    &= 2 P(\mathcal{N}(\mu_{Y_{MI}[k]}, \sigma^2_{Y_{MI}[k]}) > 0) - 1 \nonumber\\
    &= 2 P(\mathcal{N}(0, \sigma^2_{Y_{MI}[k]}) \leq \mu_{Y_{MI}[k]}) - 1,
\end{align}
%
which can be expressed in terms of the Gauss error function:
%
\begin{align}
    \alpha[k] 
    &= \mathrm{erf}\left(\frac{\mu_{Y_{MI}[k]}}{\sqrt{2 \sigma^2_{Y_{MI}[k]}}}\right) \nonumber \\
    &= \mathrm{erf}\left(
        \frac{|\mu_{Y_M[k-1]} - \mu_{Y_I[k]}|}
             {\sqrt{2[\sigma^2_{Y_M[k-1]} + \sigma^2_{Y_I[k]} - 2 \mathrm{Cov}(Y_M[k-1], Y_I[k])]}}
        \right).
\label{eq:alpha4}
\end{align}

Applying \cref{assumption:independence}, \cref{assumption:variances} and \cref{assumption:means}, 
\cref{eq:alpha4} becomes
%
\begin{equation}
    \alpha[k] 
    = \mathrm{erf}\left(
        \frac{|y_M[k-1] - y_I[k]|}
             {\sqrt{2\sigma^2_{Y_I[k]}}}
        \right)
    = \mathrm{erf}\left(
        \frac{|y_M[k-1] - y_I[k]|}
             {\sqrt{2\sigma^2_n[k]}}
        \right).
\end{equation}

Since knowledge of the measurement noise variance $\sigma^2_n[k]$ is not always available in practice,
a fixed $\lambda \geq 1$ is introduced to scale the measurement noise standard deviation.
Thus, the final expression for $\alpha[k]$ is given by
%
\begin{equation}
    \alpha[k] 
    = \mathrm{erf}\left(
        \frac{|y_M[k-1] - y_I[k]|}
             {\lambda\sqrt{2\sigma^2_n[k]}}
        \right)
\end{equation}

which is equivalent to \cref{eq:alpha_lambda} in \cref{chap:theory-framework}.